\section{The first turning point}
\label{sec:first}

At $\lambda = +\infty$ variance is irrelevant and \eqref{eq:qp} reduces to earning the
most expected return the constraints allow. For a single all-ones budget row
$\mathbf{1}\trans w = b$, \code{init\_algo}
(\srcfile{src/cvxcla/first.py}{first.py}) computes this greedily: start every weight at
$\lb$, sort the assets by descending $\mu$, and raise each towards $\ub_i$ until the
budget is met (within a tolerance of $10^{-12}$; without it, round-off can skip the
asset that should end strictly inside its box). The last asset touched absorbs the residual and is the
single \emph{free} asset; all others are blocked. The same fill serves leverage
($b > 1$) and dollar-neutral books ($b = 0$, negative $\lb$). The last asset is marked
free even when it lands exactly on a bound, so the reduced KKT system is non-singular
from the outset, a case the review of \citet{markowitz2021} does not address.

For a general equality system (weighted rows, or $m > 1$ rows such as sector or factor
neutrality) or any $G$, the maximum-return vertex is the linear program
$\max_w \mu\trans w$ subject to the constraints, which \pkg{cvxcla} solves with the
HiGHS simplex through \code{scipy.optimize.linprog}; the free set is read off as the
assets strictly inside their bounds. Only this first step differs; the loop that
follows is written for a general $A$ throughout.

\paragraph{Tied maximum returns.}
The maximum-return set need not be a single vertex. When two or more assets share the
largest expected return, say because they were given the same forecast, the linear
program has a whole face of optima, and the greedy fill or the simplex returns one of
its vertices. The frontier does not start at an arbitrary one. For $\lambda$ just below
$+\infty$ the optimiser of \eqref{eq:qp} is the point of that face where the quadratic
term $\tfrac12 w\trans\Sig w$ is smallest, and the path leaves from there. Since
release~2.6.0, \pkg{cvxcla} starts at that point, a small quadratic program over the
face, and its free set is the assets strictly inside their bounds there. Earlier
releases started at the vertex the solver returned. Release~2.3.3, used for the first
version of this paper, then traced a frontier that ended but was not efficient; that
version reported such problems as completed (Appendix~\ref{sec:degeneracy}).
Releases~2.4.2 to 2.5.0 declined them instead.

The vertex itself can be degenerate, with every weight on a bound or a duplicated tight
row. Its partition is then read from a vertex of the linear program's optimal dual set,
which keeps it dual feasible and starts the first segment at the vertex itself
(Appendix~\ref{sec:vertexdetail}). Redundant equality rows, a repeated budget say, are
dropped before the trace starts; a vertex is refused with a \code{DegenerateProblemError}
(\S\ref{sec:errors}) only when its optimal dual set has no vertex.
